\documentclass[11pt,a4paper]{ctexart}
\usepackage[margin=2.2cm]{geometry}
\usepackage{xcolor}
\usepackage{enumitem}
\usepackage{amssymb}
\usepackage[normalem]{ulem}
\usepackage{titlesec}
\usepackage{hyperref}
\definecolor{navy}{HTML}{1A365D}
\definecolor{ink}{HTML}{1F2933}
\definecolor{muted}{HTML}{66788A}
\definecolor{soft}{HTML}{EEF4FA}
\definecolor{gold}{HTML}{8A5A12}
\hypersetup{colorlinks=true,linkcolor=navy,urlcolor=navy}
\setlist{nosep}
\titleformat{\section}{\Large\bfseries\color{navy}}{}{0pt}{}
\titleformat{\subsection}{\large\bfseries\color{navy}}{}{0pt}{}
\pagestyle{plain}

\title{\textbf{七上数词与介词专项}\\[4pt]
\large 人教版（新课纲版）七年级上册 · 数词 / 介词}
\author{初中英语资料站}
\date{2026-11-12}
\begin{document}
\maketitle

\section*{使用说明}
\begin{itemize}[leftmargin=1.4em]
\item 本讲义对应人教版（新课纲版）七年级上册语法要求：数词与介词。
\item 共 4 讲语法、6 级 48 题。
\item 在线版：\url{https://1901151458-prog.github.io/english-materials/g7a-numbers/}
\end{itemize}

\section{基数词}
基数词表示数目的多少。
\begin{center}
\begin{tabular}{p{4.4cm}p{4.4cm}p{4.4cm}}
\hline
\textbf{1--10} & \textbf{11--19} & \textbf{整十}\\
\hline
one, two, three, four, five, six, seven, eight, nine, ten & eleven, twelve, thirteen, fourteen, fifteen, sixteen, seventeen, eighteen, nineteen & twenty, thirty, forty, fifty, sixty, seventy, eighty, ninety\\
\hline
\end{tabular}
\end{center}
\textbf{易错：}forty 不含 u；twenty-one 要加连字符。

\section{序数词}
序数词表示顺序的先后，\textbf{前面常用定冠词 the 或物主代词}。
\begin{center}
\begin{tabular}{p{4.6cm}p{8.4cm}}
\hline
\textbf{变化} & \textbf{例}\\
\hline
不规则 & first (1st)、second (2nd)、third (3rd)\\
拼写特殊 & fifth (5th)、eighth (8th)、ninth (9th)、twelfth (12th)\\
整十 y 变 ie 加 th & twentieth (20th)、fortieth (40th)\\
几十几只变个位 & twenty-first (21st)\\
\hline
\end{tabular}
\end{center}

\section{时间介词 at / on / in}
\begin{center}
\begin{tabular}{p{4.6cm}p{4.4cm}p{4.2cm}}
\hline
\textbf{at} & \textbf{on} & \textbf{in}\\
\hline
钟点 at 3 o'clock & 日期 on 22nd November & 月份 in April\\
年龄 at the age of 20 & 星期 on Thursday & 季节 in spring\\
at night & 节假日 on New Year's Day & 年份 in 1949\\
\hline
\end{tabular}
\end{center}

\section{地点介词与搭配}
\begin{itemize}[leftmargin=1.4em]
\item \textbf{at} 某地点：at school、at home、at the bus stop
\item \textbf{on} 表面：on the desk、on the door、on the sports field
\item \textbf{in} 范围：in China、in the garden、in Class 2
\item \textbf{方位：}under、in front of、behind、next to、across from、between
\end{itemize}
\textbf{固定搭配：}be good at、be famous for、be interested in、look at、listen to、ask for、on the way to、be similar to。\\
\textbf{注意：}介词宾语若是人称代词，要用\textbf{宾格}。


\newpage
\section{分级练习（48 题）}
\noindent 按难度分级排列，先做基础，再做成分与扩展，最后做综合题。

\subsection*{第 1 级 · 基数词}
\noindent\textcolor{gray}{1–100 的数字拼写。}\par\smallskip
\begin{enumerate}[leftmargin=*,itemsep=4pt]
\item 7 的英语是 \_\_\_。
\begin{enumerate}[label=\Alph*.,leftmargin=2em,itemsep=1pt,topsep=2pt]
\item seven
\item seventh
\item seven’s
\item seventeen
\end{enumerate}
\item 12 的英语是 \_\_\_。
\begin{enumerate}[label=\Alph*.,leftmargin=2em,itemsep=1pt,topsep=2pt]
\item twelve
\item twelveth
\item twelfth
\item twenty
\end{enumerate}
\item 20 的英语是 \_\_\_。
\begin{enumerate}[label=\Alph*.,leftmargin=2em,itemsep=1pt,topsep=2pt]
\item twenty
\item twentieth
\item twoty
\item twelve
\end{enumerate}
\item 40 的英语是 \_\_\_。
\begin{enumerate}[label=\Alph*.,leftmargin=2em,itemsep=1pt,topsep=2pt]
\item forty
\item fourty
\item fourteen
\item four
\end{enumerate}
\item 100 的英语是 \_\_\_。
\begin{enumerate}[label=\Alph*.,leftmargin=2em,itemsep=1pt,topsep=2pt]
\item one hundred
\item one hundreds
\item a hundred’s
\item hundred
\end{enumerate}
\item 99 的英语是 \_\_\_。
\begin{enumerate}[label=\Alph*.,leftmargin=2em,itemsep=1pt,topsep=2pt]
\item ninety-nine
\item ninetynine
\item ninety-ninth
\item ninty-nine
\end{enumerate}
\item 21 的英语是 \_\_\_。
\begin{enumerate}[label=\Alph*.,leftmargin=2em,itemsep=1pt,topsep=2pt]
\item twenty-one
\item twenty one
\item twentyfirst
\item twentieth-one
\end{enumerate}
\item 101 的英语是 \_\_\_。
\begin{enumerate}[label=\Alph*.,leftmargin=2em,itemsep=1pt,topsep=2pt]
\item one hundred and one
\item one hundred one
\item hundred and one
\item one hundreds and one
\end{enumerate}
\end{enumerate}

\subsection*{第 2 级 · 序数词}
\noindent\textcolor{gray}{第几的拼写与特殊变化。}\par\smallskip
\begin{enumerate}[leftmargin=*,itemsep=4pt]
\item 第 1 的英语是 \_\_\_。
\begin{enumerate}[label=\Alph*.,leftmargin=2em,itemsep=1pt,topsep=2pt]
\item first
\item one
\item firist
\item oneth
\end{enumerate}
\item 第 2 的英语是 \_\_\_。
\begin{enumerate}[label=\Alph*.,leftmargin=2em,itemsep=1pt,topsep=2pt]
\item second
\item two
\item secend
\item twoth
\end{enumerate}
\item 第 3 的英语是 \_\_\_。
\begin{enumerate}[label=\Alph*.,leftmargin=2em,itemsep=1pt,topsep=2pt]
\item third
\item three
\item thirth
\item threeth
\end{enumerate}
\item 第 5 的英语是 \_\_\_。
\begin{enumerate}[label=\Alph*.,leftmargin=2em,itemsep=1pt,topsep=2pt]
\item fifth
\item fiveth
\item five
\item fivth
\end{enumerate}
\item 第 8 的英语是 \_\_\_。
\begin{enumerate}[label=\Alph*.,leftmargin=2em,itemsep=1pt,topsep=2pt]
\item eighth
\item eightth
\item eight
\item eightst
\end{enumerate}
\item 第 9 的英语是 \_\_\_。
\begin{enumerate}[label=\Alph*.,leftmargin=2em,itemsep=1pt,topsep=2pt]
\item ninth
\item nineth
\item nine
\item ninthth
\end{enumerate}
\item 第 20 的英语是 \_\_\_。
\begin{enumerate}[label=\Alph*.,leftmargin=2em,itemsep=1pt,topsep=2pt]
\item twentieth
\item twentyth
\item twenty
\item twentith
\end{enumerate}
\item 第 21 的英语是 \_\_\_。
\begin{enumerate}[label=\Alph*.,leftmargin=2em,itemsep=1pt,topsep=2pt]
\item twenty-first
\item twenty-firsth
\item twentieth-first
\item twenty-one
\end{enumerate}
\end{enumerate}

\subsection*{第 3 级 · 时间介词 at / on / in}
\noindent\textcolor{gray}{at 钟点、on 日期、in 月份季节年份。}\par\smallskip
\begin{enumerate}[leftmargin=*,itemsep=4pt]
\item \_\_\_ 3 o’clock
\begin{enumerate}[label=\Alph*.,leftmargin=2em,itemsep=1pt,topsep=2pt]
\item at
\item on
\item in
\item 不填
\end{enumerate}
\item \_\_\_ Thursday
\begin{enumerate}[label=\Alph*.,leftmargin=2em,itemsep=1pt,topsep=2pt]
\item at
\item on
\item in
\item 不填
\end{enumerate}
\item \_\_\_ spring
\begin{enumerate}[label=\Alph*.,leftmargin=2em,itemsep=1pt,topsep=2pt]
\item at
\item on
\item in
\item 不填
\end{enumerate}
\item \_\_\_ 1949
\begin{enumerate}[label=\Alph*.,leftmargin=2em,itemsep=1pt,topsep=2pt]
\item at
\item on
\item in
\item 不填
\end{enumerate}
\item \_\_\_ New Year’s Day
\begin{enumerate}[label=\Alph*.,leftmargin=2em,itemsep=1pt,topsep=2pt]
\item at
\item on
\item in
\item 不填
\end{enumerate}
\item \_\_\_ night
\begin{enumerate}[label=\Alph*.,leftmargin=2em,itemsep=1pt,topsep=2pt]
\item at
\item on
\item in
\item 不填
\end{enumerate}
\item \_\_\_ the afternoon
\begin{enumerate}[label=\Alph*.,leftmargin=2em,itemsep=1pt,topsep=2pt]
\item at
\item on
\item in
\item 不填
\end{enumerate}
\item \_\_\_ the age of 20
\begin{enumerate}[label=\Alph*.,leftmargin=2em,itemsep=1pt,topsep=2pt]
\item at
\item on
\item in
\item 不填
\end{enumerate}
\end{enumerate}

\subsection*{第 4 级 · 地点与方位介词}
\noindent\textcolor{gray}{at/on/in 的地点用法与六个方位介词。}\par\smallskip
\begin{enumerate}[leftmargin=*,itemsep=4pt]
\item \_\_\_ school（在某地点）
\begin{enumerate}[label=\Alph*.,leftmargin=2em,itemsep=1pt,topsep=2pt]
\item at
\item on
\item in
\item 不填
\end{enumerate}
\item \_\_\_ the desk（在某表面上）
\begin{enumerate}[label=\Alph*.,leftmargin=2em,itemsep=1pt,topsep=2pt]
\item at
\item on
\item in
\item 不填
\end{enumerate}
\item \_\_\_ China（在某范围内）
\begin{enumerate}[label=\Alph*.,leftmargin=2em,itemsep=1pt,topsep=2pt]
\item at
\item on
\item in
\item 不填
\end{enumerate}
\item The shoes are \_\_\_ the bed.（在……下面）
\begin{enumerate}[label=\Alph*.,leftmargin=2em,itemsep=1pt,topsep=2pt]
\item under
\item on
\item in
\item behind
\end{enumerate}
\item The dining hall is \_\_\_ our classroom building.（在……对面）
\begin{enumerate}[label=\Alph*.,leftmargin=2em,itemsep=1pt,topsep=2pt]
\item across from
\item behind
\item next to
\item between
\end{enumerate}
\item A baby duck is \_\_\_ the big tree.（在……后面）
\begin{enumerate}[label=\Alph*.,leftmargin=2em,itemsep=1pt,topsep=2pt]
\item behind
\item in front of
\item next to
\item under
\end{enumerate}
\item This week I sit \_\_\_ my best friend, Han Lin.（紧邻）
\begin{enumerate}[label=\Alph*.,leftmargin=2em,itemsep=1pt,topsep=2pt]
\item next to
\item across from
\item between
\item behind
\end{enumerate}
\item The student centre is \_\_\_ the library and the gym.（在……中间）
\begin{enumerate}[label=\Alph*.,leftmargin=2em,itemsep=1pt,topsep=2pt]
\item between
\item among
\item in
\item next to
\end{enumerate}
\end{enumerate}

\subsection*{第 5 级 · 介词搭配}
\noindent\textcolor{gray}{be good at、look at、listen to 等。}\par\smallskip
\begin{enumerate}[leftmargin=*,itemsep=4pt]
\item be good \_\_\_ English
\begin{enumerate}[label=\Alph*.,leftmargin=2em,itemsep=1pt,topsep=2pt]
\item at
\item in
\item for
\item to
\end{enumerate}
\item be famous \_\_\_
\begin{enumerate}[label=\Alph*.,leftmargin=2em,itemsep=1pt,topsep=2pt]
\item at
\item for
\item in
\item to
\end{enumerate}
\item be interested \_\_\_
\begin{enumerate}[label=\Alph*.,leftmargin=2em,itemsep=1pt,topsep=2pt]
\item at
\item for
\item in
\item to
\end{enumerate}
\item look \_\_\_（看）
\begin{enumerate}[label=\Alph*.,leftmargin=2em,itemsep=1pt,topsep=2pt]
\item at
\item to
\item for
\item in
\end{enumerate}
\item listen \_\_\_（听）
\begin{enumerate}[label=\Alph*.,leftmargin=2em,itemsep=1pt,topsep=2pt]
\item at
\item to
\item for
\item in
\end{enumerate}
\item ask \_\_\_（请求）
\begin{enumerate}[label=\Alph*.,leftmargin=2em,itemsep=1pt,topsep=2pt]
\item at
\item to
\item for
\item in
\end{enumerate}
\item on the way \_\_\_ school
\begin{enumerate}[label=\Alph*.,leftmargin=2em,itemsep=1pt,topsep=2pt]
\item to
\item for
\item at
\item in
\end{enumerate}
\item be similar \_\_\_
\begin{enumerate}[label=\Alph*.,leftmargin=2em,itemsep=1pt,topsep=2pt]
\item at
\item to
\item in
\item for
\end{enumerate}
\end{enumerate}

\subsection*{第 6 级 · 综合提升}
\noindent\textcolor{gray}{数词与介词混合辨析。}\par\smallskip
\begin{enumerate}[leftmargin=*,itemsep=4pt]
\item 介词宾语若是人称代词，要用 \_\_\_。
\begin{enumerate}[label=\Alph*.,leftmargin=2em,itemsep=1pt,topsep=2pt]
\item 宾格
\item 主格
\item 所有格
\item 反身代词
\end{enumerate}
\item \_\_\_ Sunday morning
\begin{enumerate}[label=\Alph*.,leftmargin=2em,itemsep=1pt,topsep=2pt]
\item at
\item on
\item in
\item 不填
\end{enumerate}
\item \_\_\_ October
\begin{enumerate}[label=\Alph*.,leftmargin=2em,itemsep=1pt,topsep=2pt]
\item at
\item on
\item in
\item 不填
\end{enumerate}
\item 第 12 的英语是 \_\_\_。
\begin{enumerate}[label=\Alph*.,leftmargin=2em,itemsep=1pt,topsep=2pt]
\item twelfth
\item twelveth
\item twelve
\item twelth
\end{enumerate}
\item 第 40 的英语是 \_\_\_。
\begin{enumerate}[label=\Alph*.,leftmargin=2em,itemsep=1pt,topsep=2pt]
\item fortieth
\item fourtieth
\item forty
\item fourty
\end{enumerate}
\item \_\_\_ duty（值日）
\begin{enumerate}[label=\Alph*.,leftmargin=2em,itemsep=1pt,topsep=2pt]
\item on
\item in
\item at
\item for
\end{enumerate}
\item \_\_\_ the future（将来）
\begin{enumerate}[label=\Alph*.,leftmargin=2em,itemsep=1pt,topsep=2pt]
\item on
\item in
\item at
\item for
\end{enumerate}
\item at the music studio 中 at 表示 \_\_\_。
\begin{enumerate}[label=\Alph*.,leftmargin=2em,itemsep=1pt,topsep=2pt]
\item 在某个地点
\item 在某个表面
\item 在某个范围
\item 在……下面
\end{enumerate}
\end{enumerate}

\newpage
\section*{答案与解析}
\small
\begin{enumerate}[leftmargin=1.6em,itemsep=2pt]
\item \textbf{B1-1 A}：基数词 seven；教材基数词表。
\item \textbf{B1-2 A}：基数词 twelve。
\item \textbf{B1-3 A}：基数词 twenty。
\item \textbf{B1-4 A}：注意 forty 不含 u。
\item \textbf{B1-5 A}：教材基数词表：one hundred。
\item \textbf{B1-6 A}：整十与个位之间加连字符。
\item \textbf{B1-7 A}：教材基数词表：twenty-one。
\item \textbf{B1-8 A}：教材基数词表：one hundred and one。
\item \textbf{B2-1 A}：序数词 first (1st)；教材序数词表。
\item \textbf{B2-2 A}：序数词 second (2nd)。
\item \textbf{B2-3 A}：序数词 third (3rd)。
\item \textbf{B2-4 A}：序数词 fifth (5th)，拼写特殊。
\item \textbf{B2-5 A}：序数词 eighth (8th)，只加 h。
\item \textbf{B2-6 A}：序数词 ninth (9th)，去掉 e。
\item \textbf{B2-7 A}：整十 y 变 ie 加 th。
\item \textbf{B2-8 A}：几十几只把个位变序数；教材序数词表。
\item \textbf{B3-1 A}：at 用于钟点；教材表7。
\item \textbf{B3-2 B}：on 用于星期几；教材表7。
\item \textbf{B3-3 C}：in 用于季节；教材表7。
\item \textbf{B3-4 C}：in 用于年份；教材表7。
\item \textbf{B3-5 B}：on 用于节假日；教材表7。
\item \textbf{B3-6 A}：at night 是固定搭配；教材表7。
\item \textbf{B3-7 C}：in the afternoon；教材表7。
\item \textbf{B3-8 A}：at the age of；教材表7。
\item \textbf{B4-1 A}：at 表示在某个地点；教材。
\item \textbf{B4-2 B}：on 表示在某个表面上。
\item \textbf{B4-3 C}：in 表示在某个范围内。
\item \textbf{B4-4 A}：under 在……下面；教材例句。
\item \textbf{B4-5 A}：across from 在……对面；教材例句。
\item \textbf{B4-6 A}：behind 在……后面；教材例句。
\item \textbf{B4-7 A}：next to 紧邻；教材例句。
\item \textbf{B4-8 A}：between ... and ... 在两者之间；教材例句。
\item \textbf{B5-1 A}：be good at 擅长；教材搭配表。
\item \textbf{B5-2 B}：be famous for 因……而出名。
\item \textbf{B5-3 C}：be interested in 对……感兴趣。
\item \textbf{B5-4 A}：look at 看；教材搭配表。
\item \textbf{B5-5 B}：listen to 听；教材搭配表。
\item \textbf{B5-6 C}：ask for 请求；教材搭配表。
\item \textbf{B5-7 A}：on the way to 在去……的路上；教材搭配表。
\item \textbf{B5-8 B}：be similar to 与……相似；教材搭配表。
\item \textbf{B6-1 A}：教材：介词宾语若是人称代词，则要用宾格。
\item \textbf{B6-2 B}：某一天的某段时间用 on；教材表7。
\item \textbf{B6-3 C}：月份前用 in；教材表7。
\item \textbf{B6-4 A}：序数词 twelfth，拼写特殊。
\item \textbf{B6-5 A}：整十 y 变 ie 加 th。
\item \textbf{B6-6 A}：on duty 值日；教材搭配表。
\item \textbf{B6-7 B}：in the future 将来；教材搭配表。
\item \textbf{B6-8 A}：at 表示在某个地点；教材地点介词部分。
\end{enumerate}

\end{document}
